%==========================================================================
\section{Dynamic semantics}
\subsection*{Operand Evaluation Order}
All operands evaluate in strict left-to-right order before any dynamic
validation occurs. The requested type T is not evaluated.
\begin{enumerate}
  \item \textbf{\code{getfield}(T, e, n):} Evaluate receiver e, then field name n.
  \item \textbf{\code{setfield}(T, e, n, v):} Evaluate receiver e, field name n, then value v.
  \item \textbf{\code{invoke}(T, e, n, $e_1, \ldots, e_k$):} Evaluate receiver
    e, method name n, then positional arguments $e_1$ through $e_k$ sequentially.
  \item \textbf{\code{new\_instance}(T, n, $e_1, \ldots, e_k$):} Evaluate class
    name n, then positional arguments $e_1$ through $e_k$ sequentially.
  \item \textbf{\code{allocate}(T, n):} Evaluate class name n.
\end{enumerate}

\subsection*{Dynamic Validation Sequence}
After operand evaluation completes, the runtime validates in this order:
\begin{enumerate}
  \item \textbf{Receiver Null Check:} Verify receiver $o \neq \tm{null}$,
    including the operand of reflect. Skipped for new\_instance and allocate,
    which have no receiver.
  \item \textbf{Member / Class Existence Check:} verify that the class C, field
    f or method m exists in the dynamic metadata tables.
  \item \textbf{Type Match Check:} verify that the requested type T matches
    the target. A field's declared type must equal T, or, when T is a class,
    be T or a subclass of T; this applies to both getfield and setfield. A
    method's declared return type must equal T. For new\_instance and
    allocate, verify that the named class is T or a subclass of T.
  \item \textbf{Arity Check:} verify that the number of supplied arguments k
    equals the target method's or constructor's declared parameter count.
    Constructors are selected by arity, as with new.
  \item \textbf{Parameter Type Compatibility Check:} verify that each
    argument's static type $A_i$ (known at compile time and passed to the
    runtime) satisfies $A_i \le P_i$ for every parameter.
  \item \textbf{Value Check:} for setfield on a class-typed field, verify
    that the value is null or that its runtime class is the field's declared
    type or a subclass of it.
\end{enumerate}
If any validation step fails, execution aborts immediately with exit status
140 and writes a diagnostic of the form \code{lo\_reflect: <reason>} to
standard error. No field is written, no method body runs, and no object is
allocated before validation passes. ReflectInfo query indices start at 0, and
an out-of-range index also aborts with 140. Fields are listed inherited-first
in declaration order; methods are listed in vtable order. Calling a query on a
null ReflectInfo is an ordinary null method call (exit 102); allocation
failure keeps exit 137.

\subsection*{Memory Management, \code{ReflectInfo} Lifetimes, and Write Barriers}
\begin{itemize}
  \item \textbf{\code{ReflectInfo} Object Lifetime \& Receiver Reachability:}
    \code{ReflectInfo} instance holds a raw pointer to static class metadata
    stored in \code{.rodata}, not a managed reference to the receiver object.
    Holding a ReflectInfo therefore does not keep the receiver alive during
    garbage collection once all other references to the receiver are dropped.
  \item \textbf{new\_instance Allocation:} after validation, new\_instance
    behaves exactly like \code{new C($v_1, \ldots, v_k$)} with the
    already-evaluated arguments. It allocates through \code{lo\_alloc}, sets
    String fields to the empty string, and runs the chosen constructor,
    including any \code{this(...)} or \code{super(...)} delegation.
  \item \textbf{allocate Allocation:} after validation, allocate calls
    \code{lo\_alloc} with the class's descriptor and stores the empty string
    into each String field, including inherited ones. No constructor runs, so
    int, bool and class-typed fields keep their defaults (0, false, null);
    LO-4 already allows objects whose fields were never assigned (LO reference
    \S4.4.2).
  \item \textbf{GC Write Barrier Routing:} when setfield mutates a reference
    field, the runtime helper routes the write through
    \code{lo\_gc\_write\_barrier} after the field's value is validated.
\end{itemize}

\subsection*{Interaction with Control Flow}
Reflective operations add no new control flow; each one either completes
normally or aborts. A return inside an invoked method or a constructor returns
from that method or constructor only, as in an ordinary call. Inside a loop,
the lookup is repeated on each iteration. while and break are unchanged.

\subsection*{Worked Example}
This traces the program in Section 1 with
\code{c = new\_instance(Counter, "Counter", 3);} in place of its new
expression:
\begin{enumerate}
  \item \code{new\_instance(Counter, "Counter", 3)} evaluates
    \code{"Counter"}, then \code{3}.
    \begin{itemize}
      \item[--] Class Counter exists and Counter $\le$ Counter.
      \item[--] The implicit constructor has arity 1, and argument int $\le$ parameter int.
      \item[--] The runtime allocates o and runs the constructor, so o.count = 3 and c = o.
    \end{itemize}
  \item \code{reflect(c)}: c evaluates to o, which is not null. The runtime
    allocates q = ReflectInfo(Counter), and info = q.
  \item \code{info.class\_name()} is an ordinary call. It allocates the string
    \code{"Counter"}, which is printed.
  \item \code{setfield(int, c, "count", 7)} evaluates c $\to$ o, then
    \code{"count"}, then 7.
    \begin{itemize}
      \item[--] o is not null, and field count : int exists.
      \item[--] The requested type int equals the declared type int.
      \item[--] The runtime stores o.count = 7.
    \end{itemize}
  \item \code{invoke(int, c, "add", 5)} evaluates c $\to$ o, then
    \code{"add"}, then 5.
    \begin{itemize}
      \item[--] o is not null, and add exists and returns int, matching the requested type.
      \item[--] Arity is 1 = 1, and argument int $\le$ parameter int.
      \item[--] add runs on o with k = 5 and returns (7 + 5) = 12.
    \end{itemize}
  \item main returns 12.
\end{enumerate}

